\begin{afterword}
\summaryhead

The post begins by granting that real cults do great harm, through practices such as sleep deprivation, isolation and taking members' money. Yet people who nervously ask ``This isn't a cult, is it?'' annoy the author, who lists what is wrong with the question.

First, cults are not a natural kind like cats and dogs. Cultishness is an attractor: a set of ordinary group failures that feed each other, and any group can drift toward it, so characteristics must be judged one at a time. Second, the nervous asker wants reassurance, and a relieved ``It's not a cult'' stops the effort that every group needs. Third, the author suspects the nervousness is fear of weirdness, the fear of lonely dissent felt by Asch's subjects: people are nervous about cryonics but not about political rallies or old religions. Unusual beliefs and lofty goals are risk factors, not the disease. Fourth, cults exploit this same fear, and the wish for certainty. Fifth, the right attitude is steady caution about what a group may become. Sixth, answering the question risks giving the reassurance a cult would give. The advice is to examine a group's reasoning directly.

\responsehead

The post's first point is sound and useful: cultishness is a matter of degree, made of separate failures that can be checked one by one. So is its fourth, that conformity keeps members in a group as surely as it keeps outsiders out. Much of the rest restates posts of the previous weeks. The attractor and the entropy image are from ``Every Cause Wants to Be a Cult''; the asker who wants reassurance, the evil teacher who would exploit it, and the swordsman are from the first of the ``Two Cult Koans.'' What is new is the diagnosis of the nervousness.

That diagnosis is offered as a suspicion and then used as a fact. It opens with ``I strongly suspect'' and ``I suspect''; the next paragraph begins ``That's why''; by the paragraph after, the nervousness is simply ``the fear of believing something that will make your friends look at you really oddly.'' No evidence is added between the suspicion and the fact. The evidence offered, the comparison with political rallies, leaves out what cryonics asks of a newcomer: tens of thousands of dollars for storage in the hope that revival may one day be possible.

The post's own account of risk sits uneasily with its complaint. It grants that socially unusual belief puts a group at risk of ingroup thinking and evaporative cooling. Nervousness that tracks weirdness therefore tracks a risk factor the post itself names. The author reports being asked the question more often about powerful AIs than about probability theory, and puts this down to how weird each sounds; the post has also named a risk specific to the AI cause, the ``Super Happy Agent.'' For the complaint to hold, the askers must be treating weirdness as the disease rather than as a risk. The post shows this only by the suspicion.

The post answers a question put to its author about the author's own causes, says so, and says it may be giving the reassurance it warns against. Its closing test sits uneasily with its own warnings: ``When you know what you're trying to do and why, you'll know'' whether a group helps you. The essay had said that a lofty goal is itself a risk factor, so knowing one's goal does not show whether a group serves it.

In short: a sound point about cultishness as a matter of degree, much of it repeated from the previous weeks, and a diagnosis of other people's nervousness that begins as a suspicion, becomes a fact without new evidence, and sets aside the cost of what is being asked and the risk factors the post itself names.
\end{afterword}
